% ===================================================================
%  TABEL Nr. 14 — Die straat. — Die winkeliers.
%  Meme regime que les precedents. Tableau a UNE SEULE COLONNE, et le
%  plus peuple du livret : quatre-vingt-dix-neuf renvois, presque
%  tous des metiers.
%
%  QUATRE MOTS-OUTILS SONT EN GRAS DANS LE FAC-SIMILE : « Til » au 1,
%  « la » et « lu » au 13, « e, » au 16. L'afrikaans a l'article et
%  le pronom, il les garde tels quels — die, dit — la ou le polonais
%  et le tcheque doivent deplacer le gras sur le mot plein voisin.
%
%  ET « tintisto » (53) EST LE TEINTURIER, NON L'ETAMEUR. Le piege
%  est reel : le tcheque a lu « stano », l'etain, et a ecrit
%  « cinar ». L'ido dit « tinto », la teinture — verwer.
%
%  Trois autres faux amis du meme tableau : « mastro-veturi » (57)
%  sont les voitures de maitre, non les vehicules de demenagement ;
%  « regulierino » (89) une religieuse ; « kompana siorino » (95) une
%  dame de compagnie.
%
%  ET LE NEERLANDAIS GUETTE PARTOUT DANS CE TABLEAU DE METIERS :
%  schoorsteenvegers, blikslager, kleermakers s'ecrivent presque
%  seuls. Ce sont skoorsteenveërs, blikslaer, kleremakers.
% ===================================================================

%%K t14-apar-1 apar
\VUsaut{4.0mm}
\VUtitre{15.0pt}{60}{TABEL Nr. 14}
\VUsaut{3.0mm}

%%K t14-tit-1 sub
\VUpk{11.4pt}{40}{Die straat. --- Die winkeliers.}
\VUsaut{3.0mm}

%%K t14-01-1 p
1. --- My vriend Jan, wat op die platteland woon, het drie dae by my familie kom
deurbring. \VUgras{Tot} dan toe het hy nie die geleentheid gehad om 'n groot stad te sien
nie; daarom gebruik ons al ons dae om te wandel, en ek verduidelik aan hom wat hy nie
ken nie.

%%K t14-02-1 p
2. --- Eers het ons die \VUgras{sterrewag} \textsuperscript{(1)} gaan besoek, wat hom
baie geboei het. Toe ons deur die \VUgras{straat} \textsuperscript{(3)} teruggekom het
waar die \VUgras{Turkse bad} \textsuperscript{(2)} is, het ons 'n \VUgras{begrafnis}
\textsuperscript{(4)} teëgekom. Voor die \VUgras{lykstoet} het die \VUgras{geestelikheid}
geloop, wat die \VUgras{lykswa} voorafgegaan het waarin die \VUgras{kis} gesit is. Baie
vriende het die \VUgras{rouende} familie vergesel.

%%K t14-02-2 p
Nadat die stoet verby was, het ons die straat oorgesteek voor die groot \VUgras{bierhuis}
\textsuperscript{(5)} waarin baie \VUgras{verbruikers} \VUgras{bier} op die terras
gedrink het, beskut deur die geglaasde \VUgras{sonskerm} \textsuperscript{(6)}.

%%K t14-03-1 p
3. --- Jan was verstom oor die \VUgras{wapenskild} wat tussen die winkel van die
\VUgras{likeurmaker} \textsuperscript{(7)} en dié van die \VUgras{musiekhandelaar}
\textsuperscript{(10)} gehang het. Hy het my na die betekenis daarvan gevra; ek het
geantwoord dat dit die Engelse \VUgras{konsulaat} \textsuperscript{(8)} aandui, en ek het
hom op die \VUgras{konsul} \textsuperscript{(9)} gewys wat op die balkon is.

%%K t14-tit-2 sub
\VUpk{10.2pt}{40}{Verkenning van die winkels.}
\VUsaut{3.0mm}

%%K t14-04-1 p
4. --- Natuurlik het ons voor die uitstalling van die \VUgras{musiekinstrumente,}
\VUgras{harmoniums,} \VUgras{orrels} en \VUgras{klaviere} stilgehou; ons het na die
geïllustreerde omslae van die \VUgras{partiture} en van die \VUgras{musiekstukke} vir
klavier gekyk, en die \VUgras{lier} van verguldde hout bewonder wat as uithangbord dien.

%%K t14-05-1 p
5. --- Die stofwolke wat die \VUgras{straatveërs} \textsuperscript{(11)} en die
\VUgras{veegwa} \textsuperscript{(12)} gemaak het, het ons gedwing om vinnig verby die
mooi \VUgras{meubelstelle} van die \VUgras{meubelhandelaar} \textsuperscript{(13)} te
loop en verby die uitstalling van \VUgras{kaggels} en \VUgras{lampe} van die
\VUgras{blikslaer} \textsuperscript{(14)}.

%%K t14-06-1 p
6. --- Origens is ons op daardie plek gedwing om die sypaadjie te verlaat, want dit was
versper deur pakke wat 'n mens uit 'n \VUgras{verhuiswa} \textsuperscript{(16)} gehaal
het om hulle in 'n \VUgras{pakhuis} \textsuperscript{(15)} te sit wat pas verhuur is.

%%K t14-06-2 p
Dit het ons belet om die mooi rytuie van die \VUgras{rytuigmaker} \textsuperscript{(17)}
te sien, maar het ons nie belet om stil te hou om na die \VUgras{pakker}
\textsuperscript{(19)} te kyk wat 'n kis vir sy buurman gemaak het, die
\VUgras{handelaar in fietse en motors} \textsuperscript{(18)}. Daarna, omdat dit al
effens laat was, het ons huis toe teruggekeer.

%%K t14-tit-3 sub
\VUpk{10.2pt}{40}{Wandeling deur die stad.}
\VUsaut{3.0mm}

%%K t14-07-1 p
7. --- Die volgende dag het my moeder aan Jan voorgestel dat hy afgeneem word, daarna het
sy my opdrag gegee om hom na die \VUgras{fotograaf} \textsuperscript{(20)} te vergesel.
Ná die ete het ons die \VUgras{werkplek} van die kunstenaar binnegegaan, waar ons taamlik
lank moes wag. Om ons besig te hou het ons na die \VUgras{portrette}
\textsuperscript{(22)} gekyk wat teen die muur hang.

%%K t14-07-2 p
Toe dit ons beurt was, het die werk nie lank geduur nie, en sodra my vriend klaar was om
voor die \VUgras{fototoestel} \textsuperscript{(21)} te poseer, het ons afgegaan.

%%K t14-08-1 p
8. --- Op die eerste verdieping het ons deur die deur van die \VUgras{haarkapper}
\textsuperscript{(23)}, wat oop is, sy kneg gesien wat die hare van 'n klant geknip het.

%%K t14-08-2 p
Toe ons in die straat aangekom het, het ons voor die \VUgras{tabakwinkel}
\textsuperscript{(24)} verbygegaan. Jan was so vermaak deur die rooi \VUgras{lantern}
\textsuperscript{(25)} en die \VUgras{dik pyp} wat as \VUgras{uithangbord}
\textsuperscript{(26)} gebruik word, dat hy teen twee ou here gebots het wat gepraat het
terwyl hulle \VUgras{snuiftabak gebruik} \textsuperscript{(27)}.

%%K t14-tit-4 sub
\VUpk{10.2pt}{40}{Die banketbakkery.}
\VUsaut{3.0mm}

%%K t14-09-1 p
9. --- Die \VUgras{banknote} en \VUgras{muntstukke} uit alle lande wat in die
\VUgras{uitstalvenster} van die \VUgras{wisselaar} \textsuperscript{(29)} uitgestal was,
het ons aandag 'n paar oomblikke getrek, maar omdat dit die uur vir die versnapering was,
het ons by die \VUgras{banketbakker} \textsuperscript{(31)} ingegaan sonder om langer
stil te hou.

%%K t14-10-1 p
10. --- Toe ons al die \VUgras{lekkernye} sien wat die winkel bevat het, \VUgras{koeke,
heuningkoeke, koekies} en allerhande \VUgras{lekkers,} kon ons nie maklik kies nie.
Uiteindelik het Jan die \VUgras{roomkoekies} gekies, en ek het net 'n klein
\VUgras{kersietert} gevat, want die soetgoed \VUgras{maak my tande seer,} en ek verlang
glad nie om te dikwels die \VUgras{tandarts} \textsuperscript{(30)} te gaan besoek nie.

%%K t14-tit-5 sub
\VUpk{10.2pt}{40}{Masjienskryf.}
\VUsaut{3.0mm}

%%K t14-11-1 p
11. --- Om aan my vriend 'n \VUgras{skryfmasjien} te wys waarvan hy hoor praat het, maar
wat hy nie geken het nie, het ons die \VUgras{skryfkamer} \textsuperscript{(32)} van my
\VUgras{neef} \textsuperscript{(34)} binnegegaan. Ons het hom aangetref terwyl hy sy
briewe aan sy \VUgras{sekretaris} \textsuperscript{(35)} gedikteer het, wat
\VUgras{stenograaf} is. Skaars was 'n brief klaar of 'n \VUgras{masjienskrywer}
\textsuperscript{(33)} het dit met sy masjien oorgeskryf. Omdat ons nie die besighede van
my bloedverwant wou onderbreek nie, het ons net 'n paar oomblikke by hom gebly.

%%K t14-12-1 p
12. --- Onder die kantoor van my neef woon 'n groot \VUgras{horlosiemaker-juwelier}
\textsuperscript{(37)} wat boonop goudsmid is. Sy pragtige winkel bevat baie
\VUgras{horlosies, pendules, sakhorlosies,} \VUgras{kettinkies} en kosbare
\VUgras{juwele,} byvoorbeeld \VUgras{pêrelsnoere,} goue \VUgras{armbande, oorringe} en
\VUgras{vingerringe} versier met \VUgras{edelstene} en \VUgras{diamante}. Een van die
uitstalvensters is besonder vir die \VUgras{goudwerk} gereserveer en bevat 'n belangrike
versameling silwer \VUgras{eetgerei}.

%%K t14-13-1 p
13. --- Toe ons om die hoek van die straat gedraai het, het ek opgemerk dat 'n mens die
\VUgras{plaatjie} \textsuperscript{(36)} verwissel het wat naby die \VUgras{nommer} van
die huis sit. Ek wou my opmerking hardop sê, toe die vrolike klanke van \VUgras{die}
militêre musiek gehoor is. Ons het die regiment tegemoet begin hardloop, sonder om te
bedink dat \VUgras{dit} voor die winkels van die \VUgras{hoedemaker}
\textsuperscript{(39)} en van die \VUgras{handskoenmaker} \textsuperscript{(40)} sou
verbygaan, en dat ons net so goed geplaas sou gewees het om sy parade te sien as ons voor
die uitstalvenster van die \VUgras{optikus} \textsuperscript{(38)} gebly het, heeltemal
voorsien van \VUgras{knypbrille, lorgnette} en \VUgras{kykers}.

%%K t14-tit-6 sub
\VUpk{10.2pt}{40}{Marsparade.}
\VUsaut{3.0mm}

%%K t14-14-1 p
14. --- Voor die \VUgras{regiment} \textsuperscript{(41)} het die \VUgras{tamboermajoor}
\textsuperscript{(42)} geloop, gevolg deur die \VUgras{tamboerspelers}
\textsuperscript{(43)}, die \VUgras{horingblasers} \textsuperscript{(44)} en die hele
\VUgras{musiekkorps}. Die \VUgras{generaal} \textsuperscript{(45)}, wat met sy adjudant
op die plein was, het naby die \VUgras{waterstraal} \textsuperscript{(49)} van die
\VUgras{fontein} \textsuperscript{(48)} stilgehou om die \VUgras{parade} by te woon.

%%K t14-14-2 p
\VUgras{Die stafoffisiere} \textsuperscript{(46)}, die \VUgras{kolonel} en die
\VUgras{luitenant-kolonel} het hom met die degen gegroet. \VUgras{Die majoors} was voor
hulle \VUgras{bataljons} \textsuperscript{(47)} en elke \VUgras{kaptein} het sy
\VUgras{kompanie} aangevoer, terwyl die \VUgras{luitenante,} \VUgras{onderluitenante} en
\VUgras{onderoffisiere} hulle gewone plek by hulle manne ingeneem het.

%%K t14-15-1 p
15. --- Baie verbygangers het hulle op \VUgras{die kruispad} \textsuperscript{(50)}
versamel; die winkeliers, die \VUgras{sambreelmaker} \textsuperscript{(51)}, die
\VUgras{verwer} \textsuperscript{(53)}, die \VUgras{wapensmid} \textsuperscript{(54)} het
uitgekom om die \VUgras{soldate} te sien. Al die voertuie, \VUgras{huurrytuie}
\textsuperscript{(52)}, \VUgras{herewaens} \textsuperscript{(57)} en ook die
\VUgras{sproeiwa} \textsuperscript{(56)}, het hulle langs die sypaadjie gestel.

%%K t14-tit-7 sub
\VUpk{10.2pt}{40}{Tonele langs die straat.}
\VUsaut{3.0mm}

%%K t14-15-2 p
'n \VUgras{Handelsreisiger} \textsuperscript{(55)} wat sy \VUgras{monsterkoffers} in die
hand gedra het, het vinnig die straat oorgesteek, en die mense wat op die
\VUgras{imperiaal} \textsuperscript{(59)} van die \VUgras{omnibus} \textsuperscript{(58)}
gesit het, het hulle omgedraai om die skouspel te geniet. 'n Arme \VUgras{straatsanger}
\textsuperscript{(60)} met sy \VUgras{draaiorrel} \textsuperscript{(61)} moes sy
\VUgras{liedjie} onderbreek, want die lawaai van die tamboere het sy stem oorstem.

%%K t14-16-1 p
16. --- Toe die regiment voor die \VUgras{stofwinkel} \textsuperscript{(64)} aangekom
het, het die \VUgras{naaisters} \textsuperscript{(63)} en die \VUgras{kleremakers}
\textsuperscript{(62)} hulle \VUgras{naaimasjiene} en hulle werk verlaat om deur die
venster te kyk. \VUgras{Die winkelier} \textsuperscript{(65)}, wat pas nuwe
\VUgras{pakke} \textsuperscript{(66)} in sy \VUgras{uitstalling} gesit het, moes die
\VUgras{lakens} \textsuperscript{(67)} en die \VUgras{linne} binnetoe laat sit wat op die
sypaadjie gestaan het, naby die \VUgras{rioolopening} \textsuperscript{(68)}, uit vrees
dat die skare hulle sou omgooi; selfs 'n ou \VUgras{smous} \textsuperscript{(69)}, by wie
'n portierster kouse afgeding het, is gedwing om 'n gang binne te gaan om sy verkoop
klaar te maak.

%%K t14-16-2 p
Ons het die regiment tot by die kaserne vergesel \VUgras{en,} baie tevrede oor ons dag,
het ons op die uur van die aandete teruggekeer.

%%K t14-tit-8 sub
\VUpk{10.2pt}{40}{Verskillende bedrywe en beroepe.}
\VUsaut{3.0mm}

%%K t14-17-1 p
17. --- Die volgende dag het my vader ons voorgestel om die drukkery van die plaaslike
koerant te besoek. Ons het dit geesdriftig aanvaar.

%%K t14-17-2 p
Die \VUgras{drukker} is ook \VUgras{boekhandelaar} \textsuperscript{(70)}
\VUgras{(=boekverkoper),} \VUgras{uitgewer,} \VUgras{papierhandelaar} en
\VUgras{boekbinder}. Hy het op die eerste verdieping die \VUgras{redaksie}
\textsuperscript{(71)} van die koerant ingerig, en ook die \VUgras{graveerwerkplek,}
waarin die \VUgras{litograwe} \textsuperscript{(72)} en die \VUgras{kopergraveurs}
\textsuperscript{(73)} werk, ten slotte die \VUgras{setwerkplek,} waarin die
\VUgras{drukwerkers} \textsuperscript{(74)} is. Die eintlike \VUgras{drukkery}
\textsuperscript{(75)}, die \VUgras{drukperse} en die verskillende masjiene is op die
grondverdieping.

%%K t14-tit-9 sub
\VUpk{10.2pt}{40}{Jan vra baie vrae.}
\VUsaut{3.0mm}

%%K t14-18-1 p
18. --- Toe ons die drukkery uitgekom het, het Jan baie verstom stilgehou voor 'n klein
geglaasde kassie wat op 'n suil gesit is, teenoor die \VUgras{garingwinkel}
\textsuperscript{(77)}, en gevra waarvoor dit gebruik word. My vader het aan hom
verduidelik dat dit 'n \VUgras{brandmelder} \textsuperscript{(76)} is. Origens was alles
in die stad vir my vriend 'n verwondering, van die eenvoudige \VUgras{klipfontein}
\textsuperscript{(78)} en die ligte \VUgras{besteldriewiel} \textsuperscript{(80)} wat
deur 'n \VUgras{joggie} \textsuperscript{(79)} in livrei bestuur word, tot die
\VUgras{poswa} \textsuperscript{(81)}.

%%K t14-19-1 p
19. --- Terwyl my vader ons 'n paar oomblikke verlaat het om met die
\VUgras{belastinggaarder} \textsuperscript{(83)} te praat, wat stilgehou het om met 'n
\VUgras{advokaat} \textsuperscript{(82)} en 'n \VUgras{notaris} \textsuperscript{(84)} te
gesels, het ons ons vermaak deur die verkeer op die straat met die oë te volg. Die mees
uiteenlopende mense het voor ons verbygegaan: 'n \VUgras{banketbakkerskneg}
\textsuperscript{(85)}, wat in 'n mandjie 'n pragtige \VUgras{pastei} gedra het, het
agter 'n \VUgras{klerehandelaar} \textsuperscript{(86)} gekom; 'n
\VUgras{lanternaansteker} \textsuperscript{(87)} het met 'n \VUgras{gasman}
\textsuperscript{(88)} gepraat, 'n \VUgras{non} \textsuperscript{(89)} wat 'n weeskind
aan die hand gehou het, was op pad terug na haar klooster.

%%K t14-tit-10 sub
\VUpk{10.2pt}{40}{Op ons terugweg huis toe.}
\VUsaut{3.0mm}

%%K t14-20-1 p
20. --- Op die oomblik dat my vader hom weer by ons gevoeg het, het Jan hardop gelag toe
hy die vuil gesig en die vreemde flenterklere van twee \VUgras{skoorsteenveërs}
\textsuperscript{(91)} sien wat heeltemal met roet bedek was.

%%K t14-21-1 p
21. --- Ek wou by die \VUgras{blommeverkoopster} \textsuperscript{(92)} 'n plant vir my
moeder koop, maar my vader het my daarop gewys dat dit baie swaar sou wees om saam te dra
en dat dit beter sou wees om \VUgras{lemoene} aan te skaf. Ons het die
\VUgras{verkoopster} \textsuperscript{(94)} genader en toe die \VUgras{goewernante}
\textsuperscript{(93)}, wat vir haar pleegkind gekies het, haar aankoop klaar gehad het,
het ons ons laat bedien.

%%K t14-22-1 p
22. --- Toe ons huis toe teruggekeer het, het ons etlike kennisse teëgekom: 'n ou
vriendin van my familie, wat haar arm op haar \VUgras{geselskapsdame}
\textsuperscript{(95)} gesteun het, die \VUgras{ingenieur} \textsuperscript{(96)} van die
brûe en paaie, en een van my niggies met haar \VUgras{kinderoppasser}
\textsuperscript{(98)}.

%%K t14-22-2 p
Daarna het ons die \VUgras{modemaakster} \textsuperscript{(97)} van my moeder teëgekom en
twee \VUgras{monnike} \textsuperscript{(99)} wat vreemd aan die stad was, wat na ons toe
gekom het om my vader na die pad na die stasie te vra.
\VUsaut{6.0mm}
\VUfilet{22mm}
